\chapter{En ballade}
\label{chap:en-ballade}

La caravane devait rester trois jours en ville. Pour Kael et Kamatsu, cela
signifiait surtout trois jours durant lesquels personne ne s'inquiéterait de
les voir disparaître pendant des heures, du moment qu'ils étaient revenus
avant le départ.

Ils avaient passé une bonne partie de l'après-midi à parcourir les rues sans
but précis. Kamatsu s'arrêtait devant les étals qui attiraient son regard,
Kael devant ceux qui vendaient des armes, du matériel de voyage ou des objets
dont il essayait de deviner l'utilité. Ils avaient traversé un marché deux
fois, emprunté plusieurs ruelles simplement pour voir où elles menaient et
étaient désormais incapables de dire avec certitude par quel chemin ils
étaient arrivés.

Cela ne les inquiétait pas particulièrement.

— On devrait rentrer, finit tout de même par dire Kael.

Kamatsu regarda le ciel entre les toits.

— On a encore le temps.

— C'est ce que tu as dit tout à l'heure.

— Et j'avais raison.

Kael lui lança un regard.

— Tout à l'heure.

Kamatsu sourit et continua d'avancer.

Ils arrivèrent sur une petite place bordée de maisons et de commerces. Une
musique s'échappait d'une taverne dont la porte était grande ouverte. Kamatsu
ralentit aussitôt, tandis que Kael continua sur quelques pas avant de
s'apercevoir qu'il ne le suivait plus.

— Non.

— Je n'ai rien dit.

— Justement.

Kamatsu écoutait déjà. La mélodie était rapide, entraînante, portée par un
instrument qu'il ne reconnut pas immédiatement. Il entendait des cordes, mais
aucune des sonorités auxquelles Thom l'avait habitué. Les notes montaient et
descendaient à une vitesse qui lui semblait presque impossible.

— Juste une chanson.

— Tu as dit qu'on devait rentrer.

— C'est toi qui l'as dit.

Kael soupira.

Quelques instants plus tard, ils entraient dans la taverne.

\scenebreak

La salle était déjà bien remplie. Des habitués occupaient le comptoir, des
voyageurs s'étaient serrés autour des tables et quelques clients restaient
debout, faute de place. Les conversations se mêlaient aux rires et aux chocs
des chopes, mais la plupart des regards revenaient régulièrement vers la
musicienne.

Elle se tenait au milieu de la salle, un étrange instrument calé entre le
menton et l'épaule. Une main courait sur quatre cordes tandis que l'autre
faisait danser un archet au-dessus d'elles. Mais ce qui surprit le plus
Kamatsu, c'était qu'elle ne restait jamais en place. Elle tournait au rythme
de la mélodie, reculait entre les tables, revenait d'un pas vif et frappait
parfois le sol du talon sans jamais interrompre son jeu.

Kael et Kamatsu trouvèrent un coin libre près d'un mur, d'où Kamatsu continua
d'observer la musicienne.

— C'est quoi ?

Kael haussa les épaules.

— Aucune idée.

La musicienne accéléra. Son archet courait désormais sur les cordes avec une
vitesse qui semblait impossible, et ses pas suivaient la même cadence.
Kamatsu marquait déjà le rythme du pied.

Puis quelques mots lui vinrent.

Il les murmura d'abord presque pour lui-même, cherchant à les faire entrer
dans la mélodie. Une phrase en appela une autre. Il recommença, changea
quelques mots, en trouva de nouveaux.

Kael tourna la tête vers lui.

— Tu fais quoi ?

Kamatsu continua encore quelques instants avant de répondre.

— Je ne sais pas encore.

Kael eut un léger sourire.

— Ça promet.

Kamatsu l'ignora. Il avait cessé de regarder la salle. Toute son attention
était tournée vers la musicienne et vers cette mélodie dont il essayait
maintenant de suivre les détours avec ses propres mots.

Sa voix monta peu à peu.

Quelques clients proches se retournèrent. La musicienne l'entendit elle aussi.
Son regard trouva celui de Kamatsu tandis qu'elle continuait de jouer. Elle
sourit.

Kamatsu sourit à son tour et poursuivit.

Lorsqu'elle passa de nouveau devant lui, elle accrocha son regard. Sans cesser
de jouer, elle recula de quelques pas vers le centre de la salle et lui
adressa un léger signe de tête.

Kamatsu hésita et regarda Kael. Son frère avait déjà compris et affichait un
sourire amusé.

— Vas-y.

Il y avait dans sa voix exactement ce qu'il fallait pour que Kamatsu comprenne
qu'il espérait désormais le voir hésiter.

Kamatsu quitta le coin où ils s'étaient installés et rejoignit la musicienne.
Elle lui fit une place sans interrompre sa mélodie, et il continua à chanter
près d'elle.

Les paroles n'étaient pas toujours les mêmes. Kamatsu en oubliait certaines,
en remplaçait d'autres et en inventait à mesure que la musique avançait.
Parfois, il devait laisser passer quelques notes avant de retrouver une phrase.
Cela faisait partie du jeu.

Puis Lynsey se remit à parcourir la salle et Kamatsu la suivit.

Au début, ses mouvements furent maladroits. Il dut éviter une chaise au
dernier moment et manqua de heurter un serveur qui traversait la salle avec
trois chopes. Quelques rires éclatèrent, et Kamatsu rit avec eux avant de
continuer.

Il avait appris avec Thom à écouter le rythme et à sentir le moment où une
histoire avait besoin de ralentir ou d'accélérer. Ici, le rythme ne lui
demandait pas de rester immobile. Il suivit la musicienne entre les tables,
frappa des mains avec les clients et finit par reprendre quelques-uns de ses
pas.

Lorsqu'elle tourna, il tourna. Lorsqu'elle bondit sur le côté, il essaya de
faire de même et faillit tomber. Kael éclata de rire, et Kamatsu lui adressa
un geste peu aimable sans interrompre son chant. Cela fit rire une deuxième
table.

Peu à peu, quelques passants s'arrêtèrent devant la porte, attirés par la
musique et la voix de Kamatsu, avant de se glisser à leur tour dans une salle
déjà bondée. Derrière son comptoir, le tavernier avait désormais bien plus à
faire que remplir quelques chopes.

La musicienne changea de mélodie sans prévenir. Kamatsu s'interrompit une
seconde, surpris, puis chercha le nouveau rythme. Quelques notes plus tard, il
avait déjà recommencé à chanter.

Quand le dernier accord résonna enfin, Kamatsu avait chaud, respirait trop
vite et souriait tellement qu'il en avait mal aux joues. La salle applaudit
et il salua maladroitement avant de rejoindre Kael.

Son frère le regarda quelques secondes avec le plus grand sérieux.

— C'était affreux.

Kamatsu reprit son souffle.

— Ils ont applaudi.

— Ils ont surtout vu ta danse.

— Tu n'aurais pas fait mieux.

— Je n'aurais pas dansé.

— Parce que tu n'en es pas capable.

Un sourire espiègle se dessina sur le visage de Kamatsu. Mais avant qu'il
puisse mettre Kael au défi de danser devant tout le monde, le tavernier se
fraya un chemin jusqu'à eux, une chope de bière dans chaque main.

— Pour vous. Vous m'avez amené du monde, dit-il avec un clin d'œil.

Il leur tendit les deux chopes. Kael prit aussitôt la sienne et regarda
Kamatsu.

— Finalement, tu danses très bien.

Kamatsu éclata de rire et leva sa bière vers son frère. Kael cogna légèrement
sa chope contre la sienne avant d'en boire une longue gorgée. Kamatsu l'imita,
et son défi fut momentanément oublié tandis qu'ils profitaient ensemble de
leur récompense.

\scenebreak

La musicienne s'appelait Lynsey Sterlin.

Kamatsu l'apprit un peu plus tard, lorsqu'elle s'accorda une pause. Lui et Kael
allèrent la rejoindre, leurs chopes à la main. De près, son instrument
intriguait Kamatsu encore davantage.

Lynsey l'observa approcher avec un sourire.

— Tu chantes depuis longtemps ?

— Un peu. J'ai appris avec quelqu'un qui a voyagé quelque temps avec notre
caravane.

— Un musicien ?

— Thom Merilin.

Le nom sembla aussitôt retenir l'attention de Lynsey.

— Tu connais Thom ?

— Oui. Pourquoi ?

— Parce que je le connais aussi.

Kael leva légèrement les sourcils tandis que Kamatsu souriait.

— Il vous a appris à chanter, à vous aussi ?

— Non. Mais je l'ai déjà entendu jouer. Et raconter des histoires.

Elle eut un sourire amusé.

— Vous savez qu'il s'est fait chasser d'une auberge parce qu'il refusait de
terminer une histoire ?

Kamatsu secoua la tête.

— Il racontait l'histoire d'un homme qui cherchait un trésor. Il s'est arrêté
juste avant de révéler où il se trouvait et a annoncé que la fin serait pour
le lendemain soir.

Kamatsu sourit. Cela ressemblait assez à Thom.

— Le tavernier n'a pas apprécié ?

— Pas du tout. Il l'a mis à la porte.

Kael eut un sourire.

— Et Thom est revenu le lendemain ?

— Non. Il est allé jouer dans une autre taverne.

Lynsey laissa passer un instant avant d'ajouter :

— Et tous les clients de la première sont venus l'écouter dans la seconde.

Kamatsu éclata de rire. Lynsey se pencha légèrement vers eux, comme pour leur
confier un secret.

— Entre nous, je le soupçonne surtout de ne pas avoir encore trouvé la fin de
son histoire quand il a commencé à la raconter.

Le regard de Kamatsu revint vers l'instrument.

— Comment ça s'appelle ?

Lynsey le souleva légèrement.

— Un violon.

Kamatsu répéta le mot.

— Je peux voir ?

Elle le lui tendit. Kamatsu le prit avec beaucoup plus de précautions qu'il
n'en avait montré en dansant. Le bois était léger, les quatre cordes tendues
au-dessus d'une caisse creuse. Il passa un doigt près de l'une d'elles sans la
toucher.

— Et ça peut vraiment jouer aussi vite ?

Lynsey sourit.

— Le violon, oui. Le violoniste, ça dépend.

Kamatsu observa encore l'instrument quelques instants.

— Je peux essayer ?

— Tu en as déjà joué ?

— Non.

Lynsey lui tendit l'archet.

— Alors essaie.

Kael leva les yeux de sa chope.

— Pourquoi cette réponse m'inquiète ?

Lynsey montra rapidement à Kamatsu comment tenir l'instrument et placer
l'archet. Il essaya de reproduire sa position, coinça le violon contre son
épaule et tira doucement l'archet sur les cordes.

Le son qui en sortit fit grimacer Kael.

Kamatsu s'arrêta.

— C'était quoi ?

— Un violon qui souffre, répondit Lynsey en éclatant de rire.

Kamatsu rit à son tour. Il lui rendit l'archet, puis le violon, mais continua
d'observer l'instrument avec curiosité.

— Il faut longtemps pour apprendre ?

— Ça dépend de ce que tu appelles apprendre.

— Jouer comme toi.

Lynsey eut un petit rire.

— Oui.

Kamatsu hocha lentement la tête.

— Et pour commencer ?

— Un violon. Beaucoup de travail. Et accepter de faire des bruits horribles
pendant quelque temps.

Kael regarda Kamatsu.

— Pour ça, tu as déjà commencé.

Kamatsu lui donna un coup de coude sans quitter le violon des yeux.

\scenebreak

Ils quittèrent la taverne alors que le jour commençait à décliner.

Sur le chemin du retour, Kamatsu ralentit devant chaque boutique où il
apercevait un instrument. Les violons attirèrent naturellement son attention,
mais quelques regards aux étiquettes suffirent à refroidir son enthousiasme.
La plupart coûtaient bien plus que ce qu'il possédait.

Cela ne l'empêcha pas de continuer à regarder.

Lorsqu'ils retrouvèrent enfin la caravane, le soir approchait. Ils revenaient
avec quelques pièces en moins, quelques souvenirs en plus et, surtout, un
vieil étui sous le bras de Kamatsu.

À l'intérieur reposait un violon qui avait connu des jours meilleurs. Le bois
était rayé, la caisse portait les traces d'anciennes réparations et l'étui
fermait à peine. Il était vieux, usé, mais encore capable de produire des
notes parfaitement justes.

À condition, bien sûr, que celui qui tenait l'archet sache s'en servir.

Kamatsu avait encore du travail.